\section{The end-to-end C14.3a/b join}
\label{sec:r18-c143ab-join}

\begin{theorem}
Retain the original physical stationary C1 source, $d=1$, all strict
gaps and nonnegative event weights, literal $\beta=0$ when allowed,
the original infinitesimal total-scale intersection, and the
complement of all previously admitted domains.  Each of the following
R17 checklist branches has a finite tensor-SOS certificate for every
independent projective angle pair and every real or complex state:
\[
\begin{array}{c|l}
 \mathrm{C14.3a}&L=x\zeta>0\text{ infinitesimal},\quad
          \eta>0\text{ bounded},\quad\zeta\text{ unbounded},\\
 \mathrm{C14.3b}&L,\eta>0\text{ infinitesimal},\quad
          \zeta\text{ a positive bounded unit}.
\end{array}
\]
No bound on $L/\eta$ or $x/\eta$ is imposed.  No source rate is
compared with a testing or state rate.  These are complete branch
admissions; the other four R17 C1 branches are unchanged.
\end{theorem}

\paragraph{Precise inherited entries.}
For compact file citations, let R13 denote the preserved tree
\path{inherited_r17/inherited_r16/Q222_Revision13}.
The following are the actual entries used, rather than a deduction
from a filename or a numerical check:
\begin{enumerate}
\item R13/\path{advances/c1_lsmall_eta_bounded/remote_reduction.tex},
      source bounds SB1--SB4 and infinitesimal-$\eta$ remote proof
      SB11--SB12.  Its header assumes $x,L$ infinitesimal,
      $\zeta$ noninfinitesimal and $\eta$ bounded, and admits the
      exact complement of the chart where both $u,v$ are infinitesimal.
\item R13/\path{advances/c2_strong_remote/strong_ii_remote.tex},
      exact homogeneous parent RR3 and identities RR4--RR9.
      These are the identities explicitly imported and adapted in
      SB11--SB12, not an unqualified import of the C2 source theorem.
\item \path{inherited_r17/proofs/c1_source_ends/large_zeta_unit_eta_all_tests.tex},
      the theorem and EU1--EU6, followed by its full normal allocation
      and final projective join.  This admits the whole unit-$\eta$
      part of C14.3a, through every actual-$j$ opening.
\item \path{inherited_r17/proofs/c1_small_eta/damped_gamma_source_subchart.tex},
      exact DG3, positive source deficits DG5--DG9, and all residual
      allocations preceding DG10; and
      \path{inherited_r17/proofs/matrix_mechanism/large_zeta_damped_extension.tex},
      LD3--LD13 for the corresponding opening-safe large-$\zeta$
      source scale.  The displayed source-ratio hypotheses in their
      older all-tests theorems are used only for DG10/LD14, the term
      now replaced by the exact CJ allocation.
\item \path{inherited_r17/proofs/root/core_sign_lemma.tex}, the exact half-core
      Schur formula and sign gauges; and DS3/LS4 in
      \path{inherited_r17/proofs/c1_small_eta/damped_state_extension.tex} and
      \path{inherited_r17/proofs/matrix_mechanism/large_zeta_damped_state_extension.tex},
      the direct $X/Z$-unbounded original-parent exit.
\item \path{proofs/c1_joint/completed_j_conductor.tex}, CJ1--CJ8
      and its full local reassembly, which removes the entire
      missing $\epsilon uvQ\xi$ on the existing completed frame.
\end{enumerate}

\paragraph{Local real-state coverage when eta is infinitesimal.}
Take the exact sign gauges to arrange $u,v\ge0$.  For each real
state, $WZ\le0$ is admitted by the pure-testing half-core lemma,
including literal $Z=0$ and both local angle axes.  If $WZ>0$, then
$Z\ne0$.  When $X/Z$ is unbounded, the pure-testing positive packet
DS3/LS4 and the remaining original nonnegative event packets give a
direct certificate.  Otherwise $X/Z$ is bounded and CJ applies.
The CJ proof imposes no condition on $W/Z$ or $Y/Z$ and uses no
angle inverse.  It includes all remaining local angles and states,
including the former normalized $Z/u^2$-infinitesimal wedge.

For C14.3b the source scale is $s=x+\eta$ with unit $a=1/\zeta$;
for C14.3a it is $s=L+\eta$ with $a=1/\zeta$ infinitesimal and
$x=aL$.  The exact factors in CJ7--CJ8 are infinitesimal in both
cases.  The local certificate therefore applies to the entire
infinitesimal-$\eta$ portion of both branches, including all
previously admitted bounded-source-ratio portions.  It does not
require separately assuming any older state-ratio or angle-ratio
admission.

\paragraph{Audit of the remote complement.}
Both source branches satisfy the SB header: $x,L$ are infinitesimal,
$\zeta$ is noninfinitesimal and $\eta$ bounded.  When $\eta$ is
infinitesimal, the actual gap implies $r\sim1$ and SB3 gives
$C,D$ unbounded.  The SB11 harmonic source quantities satisfy
\[
 \gamma=\frac{yCD}{C+D},\qquad
 \gamma\sim L/2,\quad A/y\sim1/(2L),\quad
 \gamma A/y\sim1/4,\quad y/\gamma\sim2\eta\in\operatorname{Inf}.
\]
These estimates are uniform through both original gap openings and
do not require either $L/\eta$ or $x/\eta$ bounded.

For noninfinitesimal $v$, the RR3--RR4 homogeneous full-parent
argument requires only that $C,D$ are unbounded; its remaining
original packets are nonnegative.  The C1 beta event is retained
as its original positive square.  This covers every $u$ and includes
$v=\infty$ and both Alice endpoints.  For noninfinitesimal $u$ and
infinitesimal or zero $v$, SB12 uses the exact RR7--RR9 decomposition
with zero alpha cap and $k=-1$.  Its harmonic normalized block has
determinant residue $1-b_0^2>0$, where
$b=1/\sqrt{1+u^2}$ and $b_0<1$.  The sole negative term has an
infinitesimal Gram allocation by the displayed SB12 identity and
$y/\gamma\in\operatorname{Inf}$.  The original alpha and beta
packets are retained.  At $u=\infty$, $b=0$ in the original
homogeneous parent, so the same proof applies.  This is the full
remote complement, not a contact-graph endpoint evaluation.

\paragraph{Source openings and the unit-eta end.}
The actual source gap is always
$j=\eta(1+xh)-1+r>0$.  Neither CJ nor the remote proof replaces it
with a positive assumption on $j_0=\eta+r-1$.  The opening-cleared
DG/LD source identities include $j_0=0$ or slightly negative values
where permitted.  SB2 gives the original $D$ numerator a positive
unit residue; hence literal $D=0$ is impossible on this source face,
rather than silently excluded.  Literal $\beta=0$ remains allowed
through $B=\beta+1/2\ge1/2$.

The positive bounded $\eta$ of C14.3a is either infinitesimal or a
positive unit.  The former case is CJ plus the remote complement
just audited.  The latter case is precisely the R17 EU theorem,
which retains every actual-$j$ and simultaneous $r,j$ opening and
includes all projective tests.  This exhausts C14.3a.  C14.3b has
$\eta$ infinitesimal by definition, so CJ plus the same remote
complement exhausts it directly.  The exact sign gauges give all
angle signs.  Applying the real-state certificates to real and
imaginary parts gives every complex state without a shared state
normalization.  Contact equality and all normals are preserved by
the exact original-parent identities.

\paragraph{Reconciliation with the original C14.3 work item.}
The R17 checklist section 4A and \path{FRONTIER_R17.json} list exactly
the two branches above.  The earlier unit-$\eta$, unit-$\zeta$
small-$L$ region is already covered by R13 RO, JO and HO, as recorded
and proved in R13/\path{reconciliation/checklist.tex}.  Their exact
entries are R13/\path{advances/revision13_c1_small_l/}
\path{r_opening_closure.tex} (RO1),
\path{j_opening_closure.tex} (JO1), and
\path{simultaneous_r_j_closure.tex} (the theorem and HO1--HO3).
For clarity, their union is exhaustive: if actual $j$ is a unit,
then $j_0=j-\eta xh$ is a positive unit and RO applies.  If actual
$j$ is infinitesimal, $\eta+r\sim1$, so $1-r$ is a unit; unit $r$
gives JO and infinitesimal $r$ gives HO.  The latter two impose no
positivity premise on $j_0$ and no comparison of $j$ with $x$.

Thus the original C14.3 acceptance criterion is met after joining
the unchanged prior admissions: the complete parent, all printed
normals, every independent rate, source opening, state and projective
test have been covered.  C14.3a and C14.3b may be marked closed and
the common missing-normal obligation removed.  This does not close
C14.1a, C14.1b, C14.2, C14.4, the remaining weak-C2 shapes, or the
global source-routing obligation U1.  No historical checker was
rerun in this join audit.
